\documentclass[12pt, a4paper]{article}

\usepackage{xcolor}
\usepackage{geometry}
\usepackage{hyperref}
\usepackage{titlesec}
\usepackage{enumitem}
\usepackage{longtable}
\usepackage{booktabs}
\usepackage{array}
\usepackage{ragged2e}
\usepackage{fancyhdr}
\usepackage{tcolorbox}
\usepackage{amssymb}
\usepackage{listings}

\geometry{margin=2.5cm}

\usepackage{xepersian}
\settextfont{Vazirmatn}
\setlatintextfont{Times New Roman}
\setdigitfont{Vazirmatn}

\hypersetup{
    colorlinks=true,
    linkcolor=blue,
    urlcolor=blue
}

\titleformat{\section}{\Large\bfseries\color{blue!60!black}}{\thesection}{1em}{}
\titleformat{\subsection}{\large\bfseries}{\thesubsection}{1em}{}
\titleformat{\subsubsection}{\normalsize\bfseries}{\thesubsubsection}{1em}{}

\setlist{nosep, leftmargin=*}

\pagestyle{fancy}
\fancyhf{}
\fancyhead[R]{\small جزوهٔ مطالعه — جلسهٔ ۶}
\fancyhead[L]{\small طراحی شیءگرا و SOLID}
\fancyfoot[C]{\thepage}
\renewcommand{\headrulewidth}{0.4pt}

\linespread{1.4}

\lstset{
    basicstyle=\ttfamily\small,
    breaklines=true,
    frame=single,
    language=Java,
    showstringspaces=false
}

\title{
    \vspace{-1cm}
    \textbf{\Huge جزوهٔ مطالعه}\\[0.3cm]
    \textbf{\LARGE جلسهٔ ۶: طراحی شیءگرا و SOLID}\\[0.5cm]
    \large مهندسی نرم‌افزار
}
\author{
    تهیه‌کننده: ابراهیم فضلی\\
    ترم: ۱۴۰۴
}
\date{\today}

\begin{document}

\maketitle
\thispagestyle{empty}

\vspace{1cm}

\begin{center}
\begin{tabular}{|l|l|}
\hline
\textbf{درس} & مهندسی نرم‌افزار \\
\hline
\textbf{جلسه} & ۶ \\
\hline
\textbf{موضوع} & طراحی شیءگرا، SOLID، Coupling/Cohesion \\
\hline
\textbf{زمان مطالعه} & ۳-۴ ساعت \\
\hline
\textbf{پیش‌نیاز} & جلسات ۱-۵ \\
\hline
\end{tabular}
\end{center}

\newpage
\tableofcontents
\newpage

% ============================================================
\section{مقدمه}

\subsection{چرا طراحی شیءگرا مهم است؟}

در جلسات قبل یاد گرفتیم:
\begin{itemize}
    \item \textbf{جلسهٔ ۱:} مهندسی نرم‌افزار چیست
    \item \textbf{جلسهٔ ۲:} نیازمندی‌ها
    \item \textbf{جلسهٔ ۳:} فرآیندهای توسعه
    \item \textbf{جلسهٔ ۴:} معماری و مدل‌سازی
    \item \textbf{جلسهٔ ۵:} Sprint Review
\end{itemize}

حالا وقت آن است که \textbf{کد بهتر} بنویسیم.

\subsection{مشکل کد بد}

تصور کنید یک کلاس \lr{Order} دارید که:
\begin{itemize}
    \item مجموع را محاسبه می‌کند
    \item در دیتابیس ذخیره می‌کند
    \item ایمیل تأیید می‌فرستد
    \item فاکتور چاپ می‌کند
    \item به انبار اطلاع می‌دهد
\end{itemize}

\textbf{مشکل چیست؟}
\begin{itemize}
    \item اگر روش ارسال ایمیل عوض شود، باید \lr{Order} را تغییر دهیم
    \item اگر دیتابیس عوض شود، باید \lr{Order} را تغییر دهیم
    \item تست کردن سخت است
    \item استفادهٔ مجدد ممکن نیست
\end{itemize}

\textbf{نتیجه:} کد بد = هزینهٔ تغییر بالا.

\subsection{راه‌حل: طراحی شیءگرا}

طراحی شیءگرا به ما کمک می‌کند:
\begin{itemize}
    \item کد را به بخش‌های کوچک تقسیم کنیم
    \item هر بخش یک مسئولیت داشته باشد
    \item تغییرات، کم‌ترین تأثیر را داشته باشند
    \item تست کردن آسان باشد
\end{itemize}

\subsection{آنچه در این جزوه یاد می‌گیرید}

\begin{itemize}
    \item چهار اصل OOP
    \item پنج اصل SOLID
    \item Coupling و Cohesion
    \item Design Patterns (مقدمه)
    \item چطور کد بد را Refactor کنیم
\end{itemize}

% ============================================================
\section{چهار اصل OOP}

\subsection{Encapsulation — کپسوله‌سازی}

\subsubsection{تعریف}
داده‌ها و متدها در یک واحد (کلاس) قرار می‌گیرند. داده‌های داخلی 
\textbf{پنهان} هستند و فقط از طریق متدهای عمومی قابل دسترسی‌اند.

\subsubsection{مثال بد}
\begin{lstlisting}
public class Product {
    public String name;
    public double price;
    public int stock;
}
\end{lstlisting}

\textbf{مشکل:} هر کسی می‌تواند \lr{price} را منفی کند.

\subsubsection{مثال خوب}
\begin{lstlisting}
public class Product {
    private String name;
    private double price;
    private int stock;
    
    public void setPrice(double price) {
        if (price < 0) {
            throw new IllegalArgumentException("Price cannot be 
negative");
        }
        this.price = price;
    }
    
    public double getPrice() {
        return price;
    }
}
\end{lstlisting}

\textbf{مزیت:} قیمت منفی ممکن نیست.

\subsection{Inheritance — ارث‌بری}

\subsubsection{تعریف}
کلاس فرزند، ویژگی‌ها و متدهای کلاس والد را به ارث می‌برد.

\subsubsection{مثال}
\begin{lstlisting}
public class User {
    protected String email;
    protected String name;
    
    public void login() { ... }
}

public class Admin extends User {
    public void manageProducts() { ... }
}
\end{lstlisting}

\subsubsection{مزیت و عیب}
\textbf{مزیت:} استفادهٔ مجدد از کد.

\textbf{عیب:} وابستگی زیاد. تغییر در \lr{User} ممکن است \lr{Admin} را 
بشکند.

\textbf{توصیه:} \textbf{ترکیب (Composition)} بهتر از \textbf{ارث‌بری 
(Inheritance)} است.

\subsection{Polymorphism — چندشکلی}

\subsubsection{تعریف}
یک نام، چند رفتار. متد مشترک، پیاده‌سازی متفاوت.

\subsubsection{مثال}
\begin{lstlisting}
public interface PaymentMethod {
    void pay(double amount);
}

public class CreditCardPayment implements PaymentMethod {
    public void pay(double amount) { ... }
}

public class PayPalPayment implements PaymentMethod {
    public void pay(double amount) { ... }
}

// استفاده
PaymentMethod payment = new CreditCardPayment();
payment.pay(100000);
\end{lstlisting}

\textbf{مزیت:} کد عمومی، رفتار خاص.

\subsection{Abstraction — انتزاع}

\subsubsection{تعریف}
پنهان کردن جزئیات پیاده‌سازی. کاربر فقط \textbf{چه کاری} انجام می‌دهد را 
می‌بیند.

\subsubsection{مثال}
\begin{lstlisting}
public interface OrderRepository {
    Order findById(int id);
    void save(Order order);
}

// کاربر فقط می‌داند save() هست، نمی‌داند چطور
\end{lstlisting}

\textbf{ابزار:} Interface، Abstract Class.

% ============================================================
\section{SOLID — پنج اصل طراحی}

\subsection{S — Single Responsibility Principle (SRP)}

\subsubsection{تعریف}
> هر کلاس فقط یک مسئولیت داشته باشد. فقط یک دلیل برای تغییر.

\subsubsection{مثال بد}
\begin{lstlisting}
public class Order {
    public void calculateTotal() { ... }
    public void save() { ... }
    public void sendEmail() { ... }
    public void printInvoice() { ... }
}
\end{lstlisting}

\textbf{مشکل:} ۴ دلیل برای تغییر.

\subsubsection{مثال خوب}
\begin{lstlisting}
public class Order {
    public void calculateTotal() { ... }
}

public class OrderRepository {
    public void save(Order order) { ... }
}

public class EmailService {
    public void sendConfirmation(Order order) { ... }
}

public class InvoiceService {
    public void print(Order order) { ... }
}
\end{lstlisting}

\textbf{مزیت:} هر کلاس یک مسئولیت.

\subsection{O — Open/Closed Principle (OCP)}

\subsubsection{تعریف}
> باز برای توسعه، بسته برای تغییر.

\subsubsection{مثال بد}
\begin{lstlisting}
public class PaymentService {
    public void pay(String method, double amount) {
        if (method.equals("card")) {
            // کارت
        } else if (method.equals("paypal")) {
            // PayPal
        } else if (method.equals("crypto")) {
            // کریپتو
        }
        // هر روش جدید = تغییر این کلاس
    }
}
\end{lstlisting}

\subsubsection{مثال خوب}
\begin{lstlisting}
public interface PaymentMethod {
    void pay(double amount);
}

public class CardPayment implements PaymentMethod { ... }
public class PayPalPayment implements PaymentMethod { ... }

// افزودن روش جدید = فقط کلاس جدید
public class CryptoPayment implements PaymentMethod { ... }
\end{lstlisting}

\textbf{مزیت:} افزودن روش جدید بدون تغییر کد قدیمی.

\subsection{L — Liskov Substitution Principle (LSP)}

\subsubsection{تعریف}
> فرزند باید بتواند جایگزین والد شود، بدون شکستن رفتار.

\subsubsection{مثال بد}
\begin{lstlisting}
public class Rectangle {
    protected int width, height;
    public void setWidth(int w) { width = w; }
    public void setHeight(int h) { height = h; }
}

public class Square extends Rectangle {
    public void setWidth(int w) {
        width = w;
        height = w; // نقض LSP
    }
}
\end{lstlisting}

\textbf{مشکل:} اگر کسی \lr{Rectangle} را با \lr{Square} جایگزین کند، رفتار 
عوض می‌شود.

\subsubsection{مثال خوب}
\begin{lstlisting}
public interface Shape {
    int area();
}

public class Rectangle implements Shape { ... }
public class Square implements Shape { ... }
\end{lstlisting}

\textbf{مزیت:} هر دو مستقل، هیچ‌کدام رفتار دیگری را نقض نمی‌کند.

\subsection{I — Interface Segregation Principle (ISP)}

\subsubsection{تعریف}
> کلاس نباید متدهای غیرضروری را پیاده کند.

\subsubsection{مثال بد}
\begin{lstlisting}
public interface PaymentGateway {
    void pay(double amount);
    void refund(double amount);
    void checkStatus(int id);
}

// کلاس ساده باید هر ۳ متد را پیاده کند
public class SimplePayment implements PaymentGateway {
    public void pay(double amount) { ... }
    public void refund(double amount) { 
        throw new UnsupportedOperationException();
    }
    public void checkStatus(int id) { 
        throw new UnsupportedOperationException();
    }
}
\end{lstlisting}

\subsubsection{مثال خوب}
\begin{lstlisting}
public interface PaymentGateway {
    void pay(double amount);
}

public interface RefundGateway {
    void refund(double amount);
}

public interface StatusGateway {
    void checkStatus(int id);
}
\end{lstlisting}

\textbf{مزیت:} هر کلاس فقط چیزی که لازم دارد را پیاده می‌کند.

\subsection{D — Dependency Inversion Principle (DIP)}

\subsubsection{تعریف}
> به انتزاع وابسته باش، نه پیاده‌سازی.

\subsubsection{مثال بد}
\begin{lstlisting}
public class OrderService {
    private PostgresOrderRepository repo;
    
    public OrderService() {
        repo = new PostgresOrderRepository(); // وابستگی مستقیم
    }
}
\end{lstlisting}

\textbf{مشکل:} تغییر دیتابیس = تغییر \lr{OrderService}.

\subsubsection{مثال خوب}
\begin{lstlisting}
public class OrderService {
    private IOrderRepository repo;
    
    public OrderService(IOrderRepository repo) {
        this.repo = repo;
    }
}

// استفاده
IOrderRepository repo = new PostgresOrderRepository();
OrderService service = new OrderService(repo);
\end{lstlisting}

\textbf{مزیت:} تغییر دیتابیس = فقط تغییر \lr{repo} در زمان ساخت.

% ============================================================
\section{SRP در HomeStore — مثال واقعی}

\subsection{کلاس بد}

\begin{lstlisting}
public class Order {
    private List<OrderItem> items;
    private Customer customer;
    
    public double calculateTotal() { ... }
    public void save() { ... }
    public void sendEmail() { ... }
    public void printInvoice() { ... }
    public void notifyWarehouse() { ... }
}
\end{lstlisting}

\textbf{مسئولیت‌ها:} محاسبه، ذخیره، ایمیل، چاپ، اطلاع.

\textbf{دلیل تغییر:} ۵ تا!

\subsection{Refactor شده}

\begin{lstlisting}
// فقط محاسبه
public class Order {
    private List<OrderItem> items;
    
    public double calculateTotal() {
        return items.stream()
            .mapToDouble(OrderItem::getSubtotal)
            .sum();
    }
}

// ذخیره
public class OrderRepository {
    public void save(Order order) { ... }
}

// ایمیل
public class EmailService {
    public void sendConfirmation(Order order) { ... }
}

// چاپ
public class InvoiceService {
    public void print(Order order) { ... }
}

// اطلاع به انبار
public class StockService {
    public void reserve(List<OrderItem> items) { ... }
}
\end{lstlisting}

\subsection{نتیجه}

\begin{itemize}
    \item \lr{Order} فقط محاسبه
    \item \lr{OrderRepository} فقط ذخیره
    \item \lr{EmailService} فقط ایمیل
    \item \lr{InvoiceService} فقط چاپ
    \item \lr{StockService} فقط انبار
\end{itemize}

\textbf{هر کلاس = یک مسئولیت.} ✅

% ============================================================
\section{Coupling و Cohesion}

\subsection{تعریف}

\begin{center}
\begin{tabular}{|l|p{9cm}|}
\hline
\textbf{مفهوم} & \textbf{توضیح} \\
\hline
Coupling & وابستگی بین کلاس‌ها — \textbf{کم} بهتر \\
\hline
Cohesion & ارتباط داخلی کلاس — \textbf{زیاد} بهتر \\
\hline
\end{tabular}
\end{center}

\subsection{Coupling کم}

\textbf{Coupling کم:} تغییر یک کلاس، بقیه را نشکند.

\textbf{مثال بد:}
\begin{lstlisting}
public class OrderService {
    public void process() {
        PostgresOrderRepository repo = new PostgresOrderRepository();
        repo.save();
    }
}
\end{lstlisting}

\textbf{Coupling بالا:} \lr{OrderService} به \lr{PostgresOrderRepository} 
وابسته است.

\textbf{مثال خوب:}
\begin{lstlisting}
public class OrderService {
    private IOrderRepository repo;
    
    public OrderService(IOrderRepository repo) {
        this.repo = repo;
    }
}
\end{lstlisting}

\textbf{Coupling پایین:} \lr{OrderService} به interface وابسته است.

\subsection{Cohesion زیاد}

\textbf{Cohesion زیاد:} همهٔ متدهای کلاس، یک هدف.

\textbf{مثال بد:}
\begin{lstlisting}
public class Utilities {
    public void sendEmail() { ... }
    public void calculateTax() { ... }
    public void formatDate() { ... }
    public void validatePhone() { ... }
}
\end{lstlisting}

\textbf{Cohesion کم:} ۴ هدف مختلف.

\textbf{مثال خوب:}
\begin{lstlisting}
public class EmailService {
    public void send() { ... }
    public void sendBulk() { ... }
    public void validate() { ... }
}

public class TaxCalculator {
    public double calculate() { ... }
    public double calculateVAT() { ... }
}
\end{lstlisting}

\textbf{Cohesion زیاد:} هر کلاس یک هدف.

% ============================================================
\section{Design Patterns — مقدمه}

\subsection{تعریف}
> راه‌حل‌های آماده برای مشکلات رایج طراحی.

\subsection{دسته‌ها}

\begin{center}
\begin{tabular}{|l|p{9cm}|}
\hline
\textbf{دسته} & \textbf{مثال} \\
\hline
Creational & Factory، Singleton، Builder \\
\hline
Structural & Adapter، Decorator، Facade \\
\hline
Behavioral & Strategy، Observer، Command \\
\hline
\end{tabular}
\end{center}

\subsection{مثال: Factory}

\begin{lstlisting}
public interface PaymentMethod {
    void pay(double amount);
}

public class PaymentFactory {
    public static PaymentMethod create(String type) {
        if (type.equals("card")) return new CardPayment();
        if (type.equals("paypal")) return new PayPalPayment();
        throw new IllegalArgumentException();
    }
}

// استفاده
PaymentMethod payment = PaymentFactory.create("card");
payment.pay(100000);
\end{lstlisting}

\textbf{مزیت:} کلاینت نمی‌داند کدام کلاس ساخته می‌شود.

\subsection{مثال: Strategy}

\begin{lstlisting}
public interface ShippingStrategy {
    double calculate(double weight);
}

public class PostShipping implements ShippingStrategy {
    public double calculate(double weight) { return weight * 5000; }
}

public class ExpressShipping implements ShippingStrategy {
    public double calculate(double weight) { return weight * 15000; }
}

// استفاده
ShippingStrategy strategy = new ExpressShipping();
double cost = strategy.calculate(2.5);
\end{lstlisting}

\textbf{مزیت:} افزودن روش جدید بدون تغییر کد.

% ============================================================
\section{سؤالات کلیدی برای مطالعه}

\subsection{سطح ۱ (پایه)}

\subsubsection{سؤال ۱: Encapsulation چیست؟}
\textbf{پاسخ:} کپسوله‌سازی داده و متدها در یک واحد (کلاس) است. داده‌های 
داخلی پنهان هستند و فقط از طریق متدهای عمومی قابل دسترسی‌اند. مزیت: تغییر 
داخلی، بدون تأثیر بیرونی.

\subsubsection{سؤال ۲: تفاوت Inheritance و Composition چیست؟}
\textbf{پاسخ:}
\begin{itemize}
    \item \textbf{Inheritance (ارث‌بری):} کلاس فرزند از والد ارث می‌برد. 
وابستگی زیاد.
    \item \textbf{Composition (ترکیب):} کلاس از کلاس دیگر استفاده می‌کند. 
انعطاف بیشتر.
\end{itemize}
توصیه: ترکیب بهتر از ارث‌بری است.

\subsubsection{سؤال ۳: SOLID چیست؟}
\textbf{پاسخ:} پنج اصل طراحی شیءگرا:
\begin{itemize}
    \item \textbf{S:} Single Responsibility
    \item \textbf{O:} Open/Closed
    \item \textbf{L:} Liskov Substitution
    \item \textbf{I:} Interface Segregation
    \item \textbf{D:} Dependency Inversion
\end{itemize}

\subsubsection{سؤال ۴: SRP چیست؟}
\textbf{پاسخ:} هر کلاس فقط یک مسئولیت داشته باشد. فقط یک دلیل برای تغییر. 
مثال: کلاس \lr{Order} نباید هم محاسبه کند، هم ذخیره، هم ایمیل بفرستد.

\subsubsection{سؤال ۵: Coupling و Cohesion چیست؟}
\textbf{پاسخ:}
\begin{itemize}
    \item \textbf{Coupling:} وابستگی بین کلاس‌ها — کم بهتر
    \item \textbf{Cohesion:} ارتباط داخلی کلاس — زیاد بهتر
\end{itemize}

\subsection{سطح ۲ (متوسط)}

\subsubsection{سؤال ۶: چطور SRP را در کد تشخیص دهیم؟}
\textbf{پاسخ:} اگر کلاس بیش از یک دلیل برای تغییر دارد، SRP نقض شده. سؤال 
بپرسید: «چند نفر ممکن است بخواهند این کلاس را تغییر دهند؟» اگر بیش از یکی، 
SRP نقض است.

\subsubsection{سؤال ۷: تفاوت OCP و DIP چیست؟}
\textbf{پاسخ:}
\begin{itemize}
    \item \textbf{OCP:} باز برای توسعه، بسته برای تغییر. افزودن ویژگی جدید 
بدون تغییر کد قدیمی.
    \item \textbf{DIP:} به انتزاع وابسته باش، نه پیاده‌سازی.
\end{itemize}
هر دو از interface استفاده می‌کنند، ولی هدف متفاوت است.

\subsubsection{سؤال ۸: چرا LSP مهم است؟}
\textbf{پاسخ:} اگر فرزند نتواند جایگزین والد شود، کدی که با والد کار 
می‌کند، با فرزند می‌شکند. این باعث باگ‌های پنهان می‌شود.

\subsubsection{سؤال ۹: چطور Coupling را کم کنیم؟}
\textbf{پاسخ:}
\begin{itemize}
    \item به interface وابسته باش، نه پیاده‌سازی
    \item Dependency Injection
    \item Event-driven
    \item لایه‌بندی درست
\end{itemize}

\subsubsection{سؤال ۱۰: Design Pattern چیست؟}
\textbf{پاسخ:} راه‌حل آماده برای مشکلات رایج طراحی. سه دسته: Creational، 
Structural، Behavioral.

\subsection{سطح ۳ (پیشرفته)}

\subsubsection{سؤال ۱۱: اگر کلاسی هم SRP را نقض کند، هم DIP، اول کدام را 
درست کنیم؟}
\textbf{پاسخ:} اول SRP. چون کلاس بزرگ را به کلاس‌های کوچک تقسیم می‌کنیم، بعد 
وابستگی‌ها ساده‌تر می‌شوند. بعد DIP را روی کلاس‌های کوچک اعمال می‌کنیم.

\subsubsection{سؤال ۱۲: چطور مطمئن شویم Refactor کد را نشکسته؟}
\textbf{پاسخ:}
\begin{itemize}
    \item قبل از Refactor تست بنویس
    \item بعد از Refactor تست‌ها را اجرا کن
    \item اگر تست‌ها پاس شدند، Refactor درست است
\end{itemize}
این اصل «Refactor با تست» است.

\subsubsection{سؤال ۱۳: آیا همیشه باید SOLID را رعایت کنیم؟}
\textbf{پاسخ:} نه. برای پروژه‌های کوچک، SOLID ممکن است پیچیدگی اضافه کند. 
SOLID برای پروژه‌هایی است که \textbf{طول عمر} دارند و \textbf{تغییر} 
می‌کنند. برای اسکریپت‌های یک‌بارمصرف، لازم نیست.

\subsubsection{سؤال ۱۴: چطور SOLID را در تیم اجرا کنیم؟}
\textbf{پاسخ:}
\begin{itemize}
    \item Code Review با چک‌لیست SOLID
    \item مثال‌های خوب در تیم
    \item آموزش
    \item ابزارهای static analysis
\end{itemize}

\subsubsection{سؤال ۱۵: اگر کد قدیمی SOLID را نقض می‌کند، چه کار کنیم؟}
\textbf{پاسخ:}
\begin{itemize}
    \item Refactor تدریجی (نه یک‌جا)
    \item هر Sprint یک کلاس
    \item تست قبل از Refactor
    \item اولویت‌بندی بر اساس تغییرات آینده
\end{itemize}

% ============================================================
\section{تمرین‌ها}

\subsection{تمرین ۱: تشخیص SRP}
کد زیر را تحلیل کنید. چند مسئولیت دارد؟
\begin{lstlisting}
public class ProductManager {
    public void addProduct(Product p) { ... }
    public void editProduct(Product p) { ... }
    public void deleteProduct(int id) { ... }
    public void searchProducts(String q) { ... }
    public void generateReport() { ... }
}
\end{lstlisting}
\textbf{راه‌حل:} ۵ مسئولیت. باید به ۵ کلاس تقسیم شود.

\subsection{تمرین ۲: Refactor با SOLID}
کد زیر را با SOLID بازنویسی کنید:
\begin{lstlisting}
public class OrderManager {
    public void createOrder(Order o) { ... }
    public void cancelOrder(int id) { ... }
    public void sendEmail(Order o) { ... }
    public void printInvoice(Order o) { ... }
    public void trackOrder(int id) { ... }
}
\end{lstlisting}
\textbf{راه‌حل:} تقسیم به \lr{OrderService}، \lr{OrderRepository}، 
\lr{EmailService}، \lr{InvoiceService}، \lr{TrackingService}.

\subsection{تمرین ۳: تشخیص Coupling}
در کد زیر، Coupling را تشخیص دهید:
\begin{lstlisting}
public class UserService {
    public void register(User u) {
        PostgresUserRepository repo = new PostgresUserRepository();
        repo.save(u);
        SmtpEmailService email = new SmtpEmailService();
        email.send(u.getEmail());
    }
}
\end{lstlisting}
\textbf{راه‌حل:} \lr{UserService} به \lr{PostgresUserRepository} و 
\lr{SmtpEmailService} وابسته است. باید به interface وابسته باشد.

\subsection{تمرین ۴: Factory}
یک \lr{PaymentFactory} بنویسید که ۳ روش پرداخت (کارت، PayPal، کریپتو) را 
پشتیبانی کند.

\subsection{تمرین ۵: Strategy}
یک \lr{ShippingStrategy} بنویسید که ۳ روش ارسال (پست، تیپاکس، پیک) را 
محاسبه کند.

% ============================================================
\section{منابع}

\subsection{منابع اصلی}
\begin{center}
\begin{tabular}{|l|l|}
\hline
\textbf{منبع} & \textbf{فصل} \\
\hline
Clean Code — Robert C. Martin & فصل ۱-۳ \\
\hline
Clean Architecture — Robert C. Martin & فصل ۷-۱۰ \\
\hline
Sommerville — Software Engineering & فصل ۷ \\
\hline
\end{tabular}
\end{center}

\subsection{منابع آنلاین}
\begin{itemize}
    \item \textbf{ویدیو:} SOLID Principles in 10 minutes
    \item \textbf{مقاله:} The S.O.L.I.D. Principles
    \item \textbf{کتاب رایگان:} Refactoring.Guru
\end{itemize}

% ============================================================
\section{چک‌لیست مطالعه}

\subsection{قبل از مطالعه}
\begin{itemize}
    \item[$\square$] جزوه را دانلود کردم
    \item[$\square$] وقت کافی دارم (۳-۴ ساعت)
\end{itemize}

\subsection{در حین مطالعه}
\begin{itemize}
    \item[$\square$] بخش ۱ (مقدمه) را خواندم
    \item[$\square$] بخش ۲ (چهار اصل OOP) را خواندم
    \item[$\square$] بخش ۳ (SOLID) را خواندم
    \item[$\square$] بخش ۴ (SRP در HomeStore) را خواندم
    \item[$\square$] بخش ۵ (Coupling/Cohesion) را خواندم
    \item[$\square$] بخش ۶ (Design Patterns) را خواندم
    \item[$\square$] بخش ۷ (سؤالات) را خواندم
    \item[$\square$] بخش ۸ (تمرین‌ها) را خواندم
\end{itemize}

\subsection{بعد از مطالعه}
\begin{itemize}
    \item[$\square$] سؤالات سطح ۱ را جواب دادم
    \item[$\square$] سؤالات سطح ۲ را جواب دادم
    \item[$\square$] سؤالات سطح ۳ را جواب دادم
    \item[$\square$] تمرین ۱ را انجام دادم
    \item[$\square$] تمرین ۲ را انجام دادم
    \item[$\square$] تمرین ۳ را انجام دادم
    \item[$\square$] تمرین ۴ را انجام دادم
    \item[$\square$] تمرین ۵ را انجام دادم
\end{itemize}

% ============================================================
\section{جمع‌بندی}

\subsection{سه نکتهٔ کلیدی}
\begin{enumerate}
    \item \textbf{SOLID = پنج اصل برای کد قابل نگهداری}
    \begin{itemize}
        \item SRP: یک مسئولیت
        \item OCP: باز برای توسعه
        \item LSP: جایگزینی
        \item ISP: interface کوچک
        \item DIP: وابسته به انتزاع
    \end{itemize}
    \item \textbf{Coupling کم، Cohesion زیاد}
    \begin{itemize}
        \item Coupling: وابستگی بین کلاس‌ها
        \item Cohesion: ارتباط داخلی کلاس
    \end{itemize}
    \item \textbf{Refactor یک عادت است}
    \begin{itemize}
        \item هر Sprint کمی
        \item با تست
        \item تدریجی
    \end{itemize}
\end{enumerate}

\subsection{سؤال نهایی}
\begin{quote}
\textbf{اگر امروز کد Sprint 1 خود را ببینید، چند تا کلاس SRP را نقض 
می‌کنند؟}
\end{quote}

\vspace{1cm}

\begin{center}
\textbf{\Large پایان جزوه}\\[0.5cm]
\textbf{\Large موفق باشید!}
\end{center}

\end{document}
